\section{A literal-assumption caveat in a current spectral theorem}
\label{app:source}

We checked Defilippis et al.~\cite{defilippis}, arXiv:2602.05846v1,
on 20 September 2026 UTC: Definitions 2.1--2.7 and 3.3,
Assumptions 2.3--3.4, Theorem 3.6, and Appendix A.2. The current record
listed only v1, dated 5 February 2026. The spectral estimator uses
\[
 M=\sum_{i=1}^n\mathcal T(Y_i)X_iX_i^\top,
 \qquad X_i\sim N(0,I_d/d),
\]
with eigenvalue-gap selection and eigenvector norm $\sqrt d$.
The preprocessing conditions are boundedness, positive essential
supremum, and a nonzero value with positive probability. We test only
the literal universal preprocessing scope of the rate claim, not other
results or a qualified theorem. This is not our main originality claim.

\begin{proposition}[Constant preprocessing fails the literal universal claim]
\label{prop:source-constant-preprocessing}
Take orthogonal target directions $w_k^*$ with $\|w_k^*\|_2^2=d$,
$1\leq k\leq m_*\leq d$, and normalized weights
\[
 a_k=\frac{k^{-\gamma}}{\left(\sum_{j=1}^{m_*}j^{-2\gamma}\right)^{1/2}},
 \qquad \gamma>\frac12.
\]
Let $g_k(z)=(z^2-1)/\sqrt2$ and
\[
 Y=\sum_{k=1}^{m_*}a_kg_k(\langle w_k^*,X\rangle)
          +\sqrt\Delta\,\xi,
 \qquad \Delta>0,\quad \xi\sim N(0,1),
\]
with $X$ and $\xi$ independent. Choose $\mathcal T(y)=1$.
For $n\geq d\geq2$, select any number $r\leq m_*$ of the ordered
eigenvectors of $M$ by a rule depending only on its eigenvalues,
normalize them to $\sqrt d$, and pad the remaining output rows by zero.
The resulting weighted matrix reconstruction error satisfies
\begin{equation}
 \sum_{k=1}^{m_*}\frac{a_k^2}{d^2}
 \mathbb E\left[
 \|\widehat w_k\widehat w_k^\top-w_k^*w_k^{*\top}\|_F^2
 \right]\geq1.
 \label{eq:source-lower-error}
\end{equation}
The expectation is over the input sample; the assertion also holds for
random orthogonal teachers independent of those inputs.
\end{proposition}
\begin{proof}
The links are even polynomials with
$\mathbb E[g_k(Z)^2]=1$ and
$\mathbb E[g_k''(Z)]=\sqrt2$ for $Z\sim N(0,1)$.
They have a uniform polynomial envelope and satisfy the displayed
nondegeneracy conditions on the target. The weights are strictly
decreasing, have sum of squares one, and satisfy the required power-law
scaling with constants independent of $m_*$, because $\gamma>1/2$.
The preprocessing is bounded, is never zero, and has the specified
positive essential supremum:
\[
 \inf\{c:\mathbb P(\mathcal T(Y)<c)=1\}
 =\inf(1,\infty)=1.
\]
Thus it meets the literal preprocessing conditions.

With this choice, $M=\sum_iX_iX_i^\top$ depends only on isotropic
Gaussian inputs. It is a rotationally invariant Wishart matrix,
independent of any information about target directions in the labels.
For $n\geq d$ its eigenvalues are distinct almost surely. Conditional on
the ordered eigenvalues, every normalized rank-one eigenprojector is
uniformly oriented. In particular, for an eigenvector $v$ normalized by
$\|v\|_2^2=d$,
\[
 \mathbb E[vv^\top\mid\operatorname{spec}(M)]=I_d,
 \qquad
 \mathbb E[(v^\top w_k^*)^2\mid\operatorname{spec}(M)]=d.
\]
These statements concern projectors and hence do not depend on how the
sign of an eigenvector is chosen. Selection of $r$ through eigenvalues
also leaves them unchanged. Consequently every selected row has
conditional normalized error
\[
 \frac{\mathbb E[\|vv^\top-w_k^*w_k^{*\top}\|_F^2
              \mid\operatorname{spec}(M)]}{d^2}
 =\frac{d^2+d^2-2d}{d^2}=2-\frac2d.
\]
A zero-padded row has normalized error one. If $p_k$ is the probability
that row $k$ is selected, the left side of
\eqref{eq:source-lower-error} is therefore
\[
 \sum_{k=1}^{m_*}a_k^2
       \left\{p_k\left(2-\frac2d\right)+(1-p_k)\right\}
 =1+\left(1-\frac2d\right)\sum_{k=1}^{m_*}a_k^2p_k\geq1.
\]
Conditioning first on a random teacher gives the same conclusion when
the teacher is independent of the inputs.
\end{proof}

The literal rate assertion combines Theorem 3.6 with the displayed
large-sample regime $\alpha=n/d\gg m_*^{2\gamma}$, in which the claimed
weighted optimal rate is of order $m_*/\alpha$. That rate can tend to
zero, whereas \eqref{eq:source-lower-error} cannot. For example, in the
subsequent large-$\alpha,m_*$ comparison, taking
$\alpha=m_*^{2\gamma+1}$ gives $m_*/\alpha=m_*^{-2\gamma}$.
The disagreement persists after the proportional high-dimensional limit,
since the lower bound holds at every finite $d\geq\max(m_*,2)$.
Likewise, normalized squared overlap of any selected row with a fixed
target direction is $1/d$, so no fixed finite $\alpha$ yields weak
recovery from this input-only estimator.

The source proof's transition from its equations (66) to (70) has a
leading coefficient containing
$\int\mathcal T(y)\mathsf Z'(y)\,\mathrm dy$.
For the constant preprocessing this is zero: $\mathsf Z$ is a Gaussian
convolution density, tends to zero at both ends, and has an integrable
derivative, so its derivative integrates to zero. Boundedness and
nonzero preprocessing do not imply a nonzero appropriately signed
coupling to the target. An added coupling condition, or a specified
valid preprocessing family, would be a substantive qualification to
check in a corrected theorem. This observation does not supply a new
learning mechanism for our project.
